\chapter{Tabela ASCII}
\label{ap:ascii}

A tabela ASCII (\emph{American Standard Code for Information Interchange}) define a correspondência entre valores numéricos inteiros e caracteres — é o que permite que um \code{char} em C seja, ao mesmo tempo, um número e um caractere. Os valores de 0 a 31 são \textbf{caracteres de controle}, não imprimíveis (como \code{\textbackslash n} e \code{\textbackslash t}, vistos no \cref{cap:entrada-saida}).

\begin{table}[h]
\centering
\begin{tabular}{@{}cl@{\hspace{1cm}}cl@{\hspace{1cm}}cl@{\hspace{1cm}}cl@{}}
\toprule
\textbf{Dec} & \textbf{Char} & \textbf{Dec} & \textbf{Char} & \textbf{Dec} & \textbf{Char} & \textbf{Dec} & \textbf{Char} \\
\midrule
32 & (espaço) & 56 & 8 & 80 & P & 104 & h \\
33 & !  & 57 & 9 & 81 & Q & 105 & i \\
34 & "  & 58 & : & 82 & R & 106 & j \\
35 & \#  & 59 & ; & 83 & S & 107 & k \\
36 & \$  & 60 & < & 84 & T & 108 & l \\
37 & \%  & 61 & = & 85 & U & 109 & m \\
38 & \&  & 62 & > & 86 & V & 110 & n \\
39 & '  & 63 & ? & 87 & W & 111 & o \\
40 & (  & 64 & @ & 88 & X & 112 & p \\
41 & )  & 65 & A & 89 & Y & 113 & q \\
42 & *  & 66 & B & 90 & Z & 114 & r \\
43 & +  & 67 & C & 91 & [ & 115 & s \\
44 & ,  & 68 & D & 92 & \textbackslash & 116 & t \\
45 & -  & 69 & E & 93 & ] & 117 & u \\
46 & .  & 70 & F & 94 & \textasciicircum & 118 & v \\
47 & /  & 71 & G & 95 & \_ & 119 & w \\
48 & 0  & 72 & H & 96 & ` & 120 & x \\
49 & 1  & 73 & I & 97 & a & 121 & y \\
50 & 2  & 74 & J & 98 & b & 122 & z \\
51 & 3  & 75 & K & 99 & c & 123 & \{ \\
52 & 4  & 76 & L & 100 & d & 124 & | \\
53 & 5  & 77 & M & 101 & e & 125 & \} \\
54 & 6  & 78 & N & 102 & f & 126 & \textasciitilde \\
55 & 7  & 79 & O & 103 & g & & \\
\bottomrule
\end{tabular}
\caption{Tabela ASCII imprimível (valores decimais 32 a 126).}
\end{table}

\begin{table}[h]
\centering
\begin{tabular}{@{}cll@{}}
\toprule
\textbf{Dec} & \textbf{Nome} & \textbf{Sequência de escape em C} \\
\midrule
0  & NUL (nulo)         & \code{\textbackslash 0} \\
7  & BEL (alarme sonoro) & \code{\textbackslash a} \\
8  & BS (backspace)      & \code{\textbackslash b} \\
9  & TAB (tabulação)     & \code{\textbackslash t} \\
10 & LF (nova linha)     & \code{\textbackslash n} \\
13 & CR (retorno de carro) & \code{\textbackslash r} \\
27 & ESC (escape)        & --- \\
\bottomrule
\end{tabular}
\caption{Caracteres de controle mais relevantes para programação em C.}
\end{table}

\begin{dica}
Como todo caractere é internamente um número, é possível fazer aritmética diretamente sobre \code{char}s. Por exemplo, \code{'a' - 'A'} vale 32 (a diferença de posição entre maiúsculas e minúsculas na tabela), e \code{('a' + 1)} vale \code{'b'} — uma técnica comum para converter maiúsculas em minúsculas ou percorrer o alfabeto em um laço.
\end{dica}
